% Compile: xelatex lettera_presentazione.tex
\documentclass[11pt,a4paper]{article}

\usepackage{fontspec}
\setmainfont{DejaVu Serif}[Scale=0.96]
\usepackage[a4paper,margin=3cm]{geometry}
\usepackage{parskip}
\usepackage{xcolor}
\usepackage[colorlinks=true,linkcolor=black,urlcolor=blue!60!black]{hyperref}
\usepackage{xspace}
\newcommand{\physis}{\textsc{Physis}\xspace}
\newcommand{\praxis}{\textsc{Praxis}\xspace}

\begin{document}

\noindent Giovanni Antonio Pezzin \\
Universit\`a di Verona \\
\href{mailto:pezzingiovanniantonio@gmail.com}{pezzingiovanniantonio@gmail.com}

\vspace{1em}
\noindent Verona, giugno 2026

\vspace{2em}
\noindent Caro Tommaso,

\vspace{0.5em}

ti scrivo per presentarti la tesi a cui lavoro e per chiederti se accetteresti di
seguirla. Si intitola \emph{Il Governo dell'Azione: Physis e Praxis nell'Et\`a della
Tecnologia}, ed \`e costruita, libro per libro, sull'impianto della \emph{Repubblica}
di Platone --- ma per rovesciarne la politica. La tesi sostiene che l'azione umana
(\praxis) trova il proprio governatore nella natura (\physis): non la natura come
collezione di oggetti, ma come principio attivo di crescita e auto-organizzazione, la
forma che ogni cosa \`e gi\`a all'opera per divenire. Recisa da questo fondamento, la
\praxis\ degenera in arbitrio della volont\`a, e la tecnica --- \emph{techne} --- da
strumento si fa gabbia.

Tre mosse vorrei rivendicare come mie. La prima \`e l'\emph{omologia}, non l'analogia:
non che l'azione retta \emph{imiti} la natura, ma che ne sia la \emph{stessa
struttura} a un diverso livello di organizzazione. La seconda \`e la \emph{prova per
costruzione}: la tesi \`e accompagnata da un motore computazionale (\physis) e da una
piattaforma (\praxis) che girano davvero, in pubblico,\footnote{\url{https://praxisweb.xyz}
e \url{https://praxisweb.xyz/physis}.} di modo che la governabilit\`a dell'azione da
parte della natura non \`e soltanto difesa ma \emph{esibita}. La terza \`e la
\emph{repubblica dal basso}: dove Platone fa discendere il bene dall'alto, dal filosofo
che lo vede, io lo faccio \emph{scoprire dal basso}, atto coerente su atto coerente,
depositato in un bene comune governato per selezione e non per comando.

Mi permetto di proportela perch\'e credo parli alla tua lingua. Il cuore della prova
non \`e speculativo ma \emph{vissuto}: ho recuperato il senso pre-platonico della
\emph{theoria} --- servizio sacro, gara, esplorazione, l'uscire del \emph{theoros}
verso l'oracolo e i giochi, non lo sguardo da poltrona --- e l'ho attraversata in
prima persona su tre territori che la contemplazione fu inventata per disprezzare: la
\textbf{carne} (da un corpo obeso alla sua forma, fenomenologia del corpo come registro
pi\`u intimo della natura), il \textbf{desiderio} (l'\emph{eros} come andare-verso, che
legge l'immagine e il gesto prima del concetto) e il \textbf{denaro} (la posta reale
come \emph{agon} e offerta, il denaro come pura relazione di valore e misura di
fiducia). Sono i tre luoghi --- corpo, immagine, denaro --- su cui il tuo lavoro mi ha
insegnato a pensare, e sospetto che proprio l\`i la tesi trovi il suo lettore pi\`u
esigente e pi\`u giusto.

Ho il testo pronto in italiano e in inglese, con un poscritto sulle macchine e
un'appendice che mostra il motore in funzione. Sarei felice di lasciartelo e di
trovarci a parlarne quando avrai tempo: vorrei capire, prima di tutto, se l'impianto
regge al tuo occhio, e dove va stretto. Ho ancora un anno davanti, e preferirei
spenderlo lavorando con te.

\vspace{1em}
\noindent Con stima e amicizia, \\
Giovanni Antonio Pezzin

\end{document}
